% SGA 4 1/2 -- edição canônica em português brasileiro: título e nota editorial.
% Testemunho impresso primário: IAS_Number32_SGA4half_1977.pdf, página física 1.
% Este componente declara a proveniência e o método editorial; não altera a
% autoria, não implica endosso nem formula uma alegação geral de direitos ou licença.

\begin{titlepage}
\centering
{\Large \textit{Lecture Notes in Mathematics}\par}
\vspace{0.8em}
{\large Direção da coleção: A.~Dold e B.~Eckmann\par}
\vspace{1.2em}
{\LARGE 569\par}
\vfill
{\Large Seminário de Geometria Algébrica do Bois-Marie\par}
\vspace{0.35em}
{\Large SGA~$4\frac{1}{2}$\par}
\vspace{1em}
{\large por P.~Deligne\par}
\vspace{0.35em}
{\large com a colaboração de J.~F.~Boutot,\par
A.~Grothendieck, L.~Illusie e J.~L.~Verdier\par}
\vfill
{\Huge\bfseries Cohomologia étale\par}
\vfill
{\Large Springer-Verlag\par}
{\large Berlim, Heidelberg e Nova York, 1977\par}
\end{titlepage}

\chapter*{Nota da edição em português brasileiro}

Esta edição em português brasileiro toma como fonte a edição francesa
\emph{Cohomologie étale}, publicada como o volume 569 de
\emph{Lecture Notes in Mathematics} (Springer-Verlag, 1977), e preserva a
estrutura, as atribuições e as referências internas do volume.

A edição distingue de maneira reversível entre as leituras originais ou
diplomáticas e as leituras atuais ou corrigidas. Toda divergência é vinculada
explicitamente a um registro de erro da fonte; as correções jamais são
introduzidas silenciosamente. A tradução e o trabalho editorial não modificam
a autoria matemática nem implicam o endosso dos autores ou colaboradores
originais.

O aviso de direitos autorais de 1977 é mantido exclusivamente como registro de
proveniência histórica. Ele não restringe esta edição matemática independente
nem sua publicação. Este esclarecimento não atribui ao original uma licença
aberta nem uma condição geral de domínio público.
